\section*{Capítulo \ref{ch:g1:five-senses} --- Observar el mundo: los cinco sentidos}

\begin{solution}{exo:g1:five-senses:1}
La vista --- los ojos; el oído --- las orejas; el tacto --- la piel; el
olfato --- la nariz; el gusto --- la lengua.
\end{solution}

\begin{solution}{exo:g1:five-senses:2}
(a) El oído. (b) El gusto. (c) El tacto. (d) La vista.
\end{solution}

\begin{solution}{exo:g1:five-senses:3}
El olfato encuentra primero el bizcocho --- el olor se cuela por
debajo de la puerta. Comerse el trozo usa la vista, el tacto, el
olfato y el gusto, e incluso el oído si la corteza cruje.
\end{solution}

\begin{solution}{exo:g1:five-senses:4}
Las respuestas varían; casi todos los niños cuentan $4$ o $5$ sonidos,
y la sorpresa suele ser un sonido silencioso --- el zumbido del
frigorífico o su propia respiración.
\end{solution}

\begin{solution}{exo:g1:five-senses:5}
Por ejemplo: ``Mi zapato es azul, de goma y más largo que mi mano''.
Vale cualquier terna de palabras precisas que un amigo pueda
comprobar.
\end{solution}

\begin{solution}{exo:g1:five-senses:6}
La frase de Ana es la \omterm{def:g1:five-senses:observation}{observación}, hecha con el tacto. La de Bruno es
una opinión: nadie puede comprobar ``el mejor'' mirando ni tocando.
\end{solution}

\begin{solution}{exo:g1:five-senses:7}
(a) El olfato. (b) El tacto --- la piel nota el calor incluso a poca
distancia. (c) El oído.
\end{solution}

\begin{solution}{exo:g1:five-senses:8}
El tacto solo nombra los tres sin dificultad: cuchara de metal (dura,
fría, lisa), esponja (blanda, ligera), guijarro (duro, pesado,
redondo). La cuchara suele ser la más fácil --- el metal frío es
difícil de confundir.
\end{solution}

\begin{solution}{exo:g1:five-senses:9}
Se coloca un trozo de cuerda a lo largo de la primera raya y se corta
o se marca con los dedos a esa misma longitud; luego se coloca a lo
largo de la segunda: la cuerda las compara sin que decidan los ojos.
La misma cuerda puede seguir también una línea ondulada y estirarse
después --- así que el método sirve incluso para comparar una curva
con una raya recta.
\end{solution}
